\section{Introduction}
This chapter presents the logical and structural design of EduRate. It describes how system components interact, how data flows across modules, and how core entities are organized in storage.

\section{Architecture Overview}
EduRate follows a client-server architecture with a web frontend, REST APIs, and a MySQL backend.

\begin{figure}[H]
\centering
\begin{tikzpicture}[
  node distance=1.8cm and 1.8cm,
  box/.style={draw, rounded corners, minimum width=3.2cm, minimum height=0.9cm, align=center, fill=blue!8},
  arr/.style={-Stealth, thick}
]
\node[box] (users) {Users\\(Student/Professor/Admin)};
\node[box, right=of users] (frontend) {Frontend\\HTML/CSS/JS};
\node[box, right=of frontend] (api) {Backend API\\Node.js + Express};
\node[box, below=of api] (db) {MySQL Database};
\node[box, below=of frontend] (email) {Email Service\\Nodemailer};

\draw[arr] (users) -- (frontend);
\draw[arr] (frontend) -- (api);
\draw[arr] (api) -- (db);
\draw[arr] (api) -- (email);
\draw[arr] (api) -- (frontend);
\end{tikzpicture}

\caption{System Architecture Diagram}
\end{figure}
\noindent\textit{This diagram presents the high-level architecture of EduRate. It shows how users interact with the frontend, backend services, and database for core workflows.}

\subsection*{Architecture Components}
\begin{enumerate}[label=\Alph*., leftmargin=1.5em]
    \item \textbf{Presentation Layer:} Browser-based UI for students, professors, and admins.
    \item \textbf{Application Layer:} Express APIs for authentication, ratings, doubts, predictions, and admin actions.
    \item \textbf{Data Layer:} MySQL schema with relational tables and analytical views.
    \item \textbf{External Service Layer:} Email communication via Nodemailer.
\end{enumerate}

\section{DFD Level 0 (Context Diagram)}
\begin{figure}[H]
\centering
\begin{tikzpicture}[
  node distance=2.1cm and 2.8cm,
  ext/.style={draw, minimum width=2.7cm, minimum height=0.9cm, fill=green!8, align=center},
  proc/.style={draw, rounded corners, minimum width=4cm, minimum height=1.1cm, fill=orange!12, align=center},
  arr/.style={-Stealth, thick}
]
\node[proc] (sys) {EduRate System};
\node[ext, left=of sys] (student) {Student};
\node[ext, above right=of sys] (prof) {Professor};
\node[ext, below right=of sys] (admin) {Admin};
\node[ext, below=of sys] (db) {MySQL DB};

\draw[arr] (student) -- (sys);
\draw[arr] (sys) -- (student);
\draw[arr] (prof) -- (sys);
\draw[arr] (sys) -- (prof);
\draw[arr] (admin) -- (sys);
\draw[arr] (sys) -- (admin);
\draw[arr] (sys) -- (db);
\draw[arr] (db) -- (sys);
\end{tikzpicture}

\caption{DFD Level 0}
\end{figure}
\noindent\textit{This context-level DFD highlights the external entities that interact with the system. It summarizes major data exchange between users and the EduRate platform.}

\subsection*{DFD Level 0 Description}
At the context level, EduRate exchanges requests and responses with three external entities (student, professor, and admin) while persisting and retrieving records from the database.

\section{DFD Level 1}
\begin{figure}[H]
\centering
\begin{tikzpicture}[
  node distance=1.2cm and 1.6cm,
  proc/.style={draw, rounded corners, minimum width=3.6cm, minimum height=0.8cm, fill=yellow!15, align=center},
  store/.style={draw, minimum width=2.4cm, minimum height=0.8cm, fill=cyan!10, align=center},
  arr/.style={-Stealth, thick}
]
\node[proc] (p1) {P1: Authentication};
\node[proc, below=of p1] (p2) {P2: Rating Management};
\node[proc, below=of p2] (p3) {P3: Doubt Management};
\node[proc, below=of p3] (p4) {P4: Performance Prediction};
\node[proc, below=of p4] (p5) {P5: Admin Analytics \& Email};

\node[store, right=2.8cm of p2] (d1) {Users};
\node[store, right=2.8cm of p3] (d2) {Ratings};
\node[store, right=2.8cm of p4] (d3) {Doubts};
\node[store, right=2.8cm of p5] (d4) {Email Logs};

\draw[arr] (p1) -- (d1);
\draw[arr] (p2) -- (d2);
\draw[arr] (p3) -- (d3);
\draw[arr] (p5) -- (d4);
\draw[arr] (p4) -- (d2);
\end{tikzpicture}

\caption{DFD Level 1}
\end{figure}
\noindent\textit{This diagram decomposes the main system into functional processes such as authentication, ratings, doubts, and reporting. It also indicates key data stores used by each process.}

\subsection*{DFD Level 1 Description}
The level-1 decomposition presents major internal processes for authentication, ratings, doubts, prediction, and administrative communication. Each process reads from or updates specific data stores.

\section{DFD Level 2 (Rating and Feedback Process)}
\begin{figure}[H]
\centering
\begin{tikzpicture}[
  node distance=1.3cm and 1.8cm,
  proc/.style={draw, rounded corners, minimum width=4.0cm, minimum height=0.9cm, fill=orange!10, align=center},
  store/.style={draw, minimum width=3.0cm, minimum height=0.9cm, fill=cyan!10, align=center},
  ext/.style={draw, minimum width=2.8cm, minimum height=0.9cm, fill=green!10, align=center},
  arr/.style={-Stealth, thick}
]
\node[ext] (student) {Student};
\node[proc, below=of student] (validate) {P2.1 Validate Rating Input};
\node[proc, below=of validate] (save) {P2.2 Store Rating and Feedback};
\node[proc, below=of save] (aggregate) {P2.3 Update Professor Aggregates};
\node[proc, below=of aggregate] (refresh) {P2.4 Return Updated Result};

\node[store, right=3.8cm of save] (ratings) {D2: Ratings Table};
\node[store, right=3.8cm of aggregate] (profs) {D3: Professors Table};

\draw[arr] (student) -- (validate);
\draw[arr] (validate) -- (save);
\draw[arr] (save) -- (aggregate);
\draw[arr] (aggregate) -- (refresh);
\draw[arr] (refresh) -- (student);

\draw[arr] (save.east) -- (ratings.west);
\draw[arr] (aggregate.east) -- (profs.west);
\draw[arr] (ratings.west) |- (aggregate.east);
\draw[arr] (profs.west) |- (refresh.east);
\end{tikzpicture}

\caption{DFD Level 2}
\end{figure}
\noindent\textit{This level-2 DFD focuses specifically on the rating and feedback pipeline. It captures validation, persistence, aggregate update, and dashboard refresh flow.}

\subsection*{DFD Level 2 Description}
The level-2 model expands the rating process into validation, storage, aggregation update, and dashboard refresh steps. This captures internal logic after student rating submission.

\section{User Journey Diagram}
\begin{figure}[H]
\centering
\begin{tikzpicture}[
  node distance=1.2cm and 1.5cm,
  step/.style={draw, rounded corners, minimum width=4.8cm, minimum height=0.9cm, fill=blue!8, align=center},
  branch/.style={draw, rounded corners, minimum width=4.8cm, minimum height=0.9cm, fill=teal!10, align=center},
  prof/.style={draw, rounded corners, minimum width=4.8cm, minimum height=0.9cm, fill=orange!12, align=center},
  arr/.style={-Stealth, thick}
]
\node[step] (reg) {Student Registration};
\node[step, below=of reg] (login) {User Login};
\node[step, below=of login] (dash) {Student Dashboard Opened};
\node[branch, below left=1.6cm and 3.2cm of dash] (rate) {Open Professor Profile and Submit Rating};
\node[branch, below right=1.6cm and 3.2cm of dash] (doubt) {Post Doubt Query for Professor};
\node[prof, below=2.2cm of doubt] (pdash) {Professor Dashboard Receives Doubt};
\node[prof, below=of pdash] (solve) {Professor Submits Solution/Answer};

\draw[arr] (reg) -- (login);
\draw[arr] (login) -- (dash);
\draw[arr] (dash) -- (rate);
\draw[arr] (dash) -- (doubt);
\draw[arr] (doubt) -- (pdash);
\draw[arr] (pdash) -- (solve);
\end{tikzpicture}

\caption{End-to-End User Journey from Registration to Doubt Resolution}
\end{figure}
\noindent\textit{This flowchart explains the complete operational journey in the application. It begins with registration and login, then shows dashboard usage, rating submission, doubt posting, and professor-side doubt resolution.}

\section{ER Diagram}
\begin{figure}[H]
\centering
\begin{tikzpicture}[
  node distance=1.7cm and 2.3cm,
  ent/.style={draw, minimum width=2.8cm, minimum height=0.9cm, fill=purple!10, align=center},
  rel/.style={-Stealth, thick}
]
\node[ent] (users) {USERS};
\node[ent, right=of users] (ratings) {RATINGS};
\node[ent, right=of ratings] (profs) {PROFESSORS};
\node[ent, below=of ratings] (doubts) {DOUBTS};
\node[ent, below=of users] (sessions) {SESSIONS};

\draw[rel] (users) -- node[above]{gives} (ratings);
\draw[rel] (ratings) -- node[above]{for} (profs);
\draw[rel] (users) -- node[left]{posts} (doubts);
\draw[rel] (profs) -- node[right]{answers} (doubts);
\draw[rel] (users) -- node[left]{has} (sessions);
\end{tikzpicture}

\caption{Entity Relationship Diagram}
\end{figure}
\noindent\textit{This ER diagram defines the structural relationships among core entities in the EduRate database. It clarifies how user, professor, rating, doubt, and session data are linked.}

\subsection*{ER Design Notes}
The schema ensures one-to-many mappings between users and ratings, professors and ratings, and users and sessions. Doubt records maintain linkage between student queries and professor responses.

\section{Use Case Diagram}
\begin{figure}[H]
\centering
\begin{tikzpicture}[
  node distance=1.2cm and 1.8cm,
  actor/.style={draw, minimum width=2.6cm, minimum height=0.8cm, fill=gray!10, align=center},
  use/.style={draw, ellipse, minimum width=4.0cm, minimum height=0.9cm, fill=teal!10, align=center},
  arr/.style={-Stealth, thick}
]
\node[actor] (student) {Student};
\node[actor, below=of student] (prof) {Professor};
\node[actor, below=of prof] (admin) {Admin};

\node[use, right=3.2cm of student] (u1) {Register / Login};
\node[use, below=of u1] (u2) {Rate Professor};
\node[use, below=of u2] (u3) {Post Doubt};
\node[use, below=of u3] (u4) {Use Performance Predictor};
\node[use, below=of u4] (u5) {View Feedback / Answer Doubts};
\node[use, below=of u5] (u6) {View Analytics / Send Emails};

\draw[arr] (student) -- (u1);
\draw[arr] (student) -- (u2);
\draw[arr] (student) -- (u3);
\draw[arr] (student) -- (u4);
\draw[arr] (prof) -- (u5);
\draw[arr] (admin) -- (u6);
\end{tikzpicture}

\caption{Use Case Diagram}
\end{figure}
\noindent\textit{This use-case view maps responsibilities of students, professors, and administrators. It shows the primary operations available to each role in the application.}

\section{Sequence Diagram (Rating Submission)}
\begin{figure}[H]
\centering
\begin{tikzpicture}[
  lifeline/.style={draw, minimum width=2.8cm, minimum height=0.9cm, fill=red!6, align=center},
  arr/.style={-Stealth, thick}
]
\node[lifeline] (s) at (0,0) {Student};
\node[lifeline] (f) at (4.8,0) {Frontend};
\node[lifeline] (a) at (9.6,0) {API};
\node[lifeline] (d) at (14.4,0) {Database};

\draw[arr] (s.east) -- node[above]{Submit rating} (f.west);
\draw[arr] (f.east) -- node[above]{API request} (a.west);
\draw[arr] (a.east) -- node[above]{Store record} (d.west);
\draw[arr] (d.west) to[bend left=12] node[below]{Write success} (a.east);
\draw[arr] (a.west) to[bend left=12] node[below]{Return updated data} (f.east);
\draw[arr] (f.west) to[bend left=12] node[below]{Show success} (s.east);
\end{tikzpicture}

\caption{Sequence Diagram for Rating Flow}
\end{figure}
\noindent\textit{This sequence diagram illustrates runtime interaction for rating submission. It shows request flow from student interface to backend and database, followed by response propagation.}

\section{Database Design Summary}
Main tables in the schema are:
\begin{enumerate}[label=\Alph*., leftmargin=1.5em]
    \item \texttt{users}
    \item \texttt{professors}
    \item \texttt{ratings}
    \item \texttt{doubts}
    \item \texttt{gifts}
    \item \texttt{sessions}
\end{enumerate}

The design supports normalized storage, referential integrity, and efficient aggregation for analytics views such as top performers and average ratings.
\noindent\textit{The table structure is designed so that each transaction can be traced from user action to persistent storage. This enables both operational reliability and analytics consistency.}

\section*{Chapter Summary}
This chapter presented the architectural and process-level design of EduRate using multiple formal diagrams. It covered DFD levels, ER modeling, sequence logic, and end-to-end user journey so that the complete data and control flow of the system remains clear and auditable.
